\section{Super-Resolution for Remote Sensing}
Traditional SR methods (bicubic, Lanczos) are fast but cannot recover high-frequency detail. Deep learning SR -- EDSR~\cite{lim2017enhanced}, SwinIR~\cite{liang2021swinir}, Real-ESRGAN~\cite{wang2021realesrgan} -- learns a mapping from LR to HR using convolutional or transformer architectures, achieving state-of-the-art fidelity on natural images.

For remote sensing, SR has been applied to Sentinel-2, Landsat, and commercial imagery. However, most evaluations rely solely on PSNR/SSIM, which do not capture spectral fidelity (critical for vegetation indices like NDVI) or structural plausibility.

\section{Hallucination in Super-Resolution}
Hallucination refers to generated details that are plausible but not grounded in the input. In satellite imagery, this includes:
\begin{itemize}
    \item \textbf{Spectral:} implausible band correlations, NDVI out of physical range
    \item \textbf{Structural:} spurious high-frequency content in homogeneous regions
    \item \textbf{Textural:} mismatched GLCM entropy / texture statistics
    \item \textbf{Semantic:} task-performance degradation vs HR reference
\end{itemize}
Existing work (e.g., Blau \& Michaeli~\cite{blau2018perception}, Ledig et al.~\cite{ledig2017photo}) discusses the perception-distortion tradeoff, but systematic hallucination detection for satellite SR remains underexplored.

\section{Task-Driven Evaluation}
Fixed-backbone evaluation (train a task model on HR, freeze, evaluate on SR outputs) provides a task-level metric normalized to HR performance. This directly measures whether SR preserves information useful for downstream tasks, beyond pixel fidelity.
